\begin{lemma}
\label{lem:p213-triv-interval-current}
Let $I$ be a connected boundaryless interval with either of the tropical
structures in \cite[\S\S6.7.1--6.7.2]{CLD25}.  Every positive
$\mathrm d'$- and $\mathrm d''$-closed current $T$ of bidegree
$(0,0)$ on $I$ has the form
\begin{equation}
                         T=c\,\delta_I,
            \qquad c\in\mathbf R_{\geq 0}.
\label{eq:p213-triv-interval-classification}
\end{equation}
Consequently $T$ is a Bedford--Taylor current in the full sense of
\cite[Definition~8.6.2]{CLD25}.
\end{lemma}

\begin{proposition}
\label{prop:p213-triv-abstract-cross-counterexample}
Let
\[
 \Sigma=\mathbf R e_1\cup\mathbf R e_2\subset\mathbf R^2
\]
be the coordinate cross, equipped with its standard tropical structure and
unit calibration on each of its four rays.  Write
$\rho_{+1},\rho_{-1},\rho_{+2},\rho_{-2}$ for the rays in the indicated
coordinate directions and set
\begin{equation}
 T=2\delta_{\rho_{+1}}+2\delta_{\rho_{-1}}
      +\delta_{\rho_{+2}}+\delta_{\rho_{-2}}.
\label{eq:p213-triv-abstract-cross-current}
\end{equation}
Then $T$ is a positive $\mathrm d'$- and $\mathrm d''$-closed current of
bidegree $(0,0)$, but it is not a Bedford--Taylor current.  Indeed, for
the continuous psh functions
\begin{equation}
                U(x,y)=2|x|,
       \qquad   V(x,y)=|x|+|y|,
\label{eq:p213-triv-abstract-cross-potentials}
\end{equation}
one has
\begin{equation}
 \mathrm d'\mathrm d''[U]=\mathrm d'\mathrm d''[V]=4\delta_0,
 \qquad
 \mathrm d'\mathrm d''(UT)=8\delta_0,
 \qquad
 \mathrm d'\mathrm d''(VT)=6\delta_0.
\label{eq:p213-triv-abstract-cross-defect}
\end{equation}
Thus the general Bedford--Taylor independence question following
\cite[\S8.6.1]{CLD25} has a negative answer for abstract tropical spaces,
already on an underlying metric tree.  This proposition makes no assertion
that $T$ lifts to a current on a Berkovich curve having $\Sigma$ as a
coarse tropical image.
\end{proposition}

\begin{proposition}
\label{prop:p213-triv-p1-current-rigidity}
Let $k$ be an algebraically closed field with the trivial absolute value,
let
\[
                         X=(\mathbf P^1_k)^{\mathrm{an}},
\]
and let $Y\subset X$ be a connected analytic open subset.  Every positive
$\mathrm d'$- and $\mathrm d''$-closed Chambert--Loir--Ducros current
$T$ of bidegree $(0,0)$ on $Y$ is a nonnegative constant multiple of
the integration current:
\begin{equation}
                         T=c\,\delta_Y,
            \qquad c\in\mathbf R_{\geq0}.
\label{eq:p213-triv-p1-classification}
\end{equation}
In particular, $T$ is a Bedford--Taylor current for arbitrary CLD-psh
pairs, not merely for continuous pairs.  This applies in particular to the
trivially valued affine line, the punctured projective line, open discs,
and punctured open discs whenever they are viewed as connected analytic
opens in $X$.
\end{proposition}

\begin{proposition}
\label{prop:p213-triv-algebraic-curve-current-rigidity}
Let $k$ be an algebraically closed field with the trivial absolute value
and let $C$ be a smooth connected projective $k$-curve.  Then every
positive $\mathrm d'$- and $\mathrm d''$-closed CLD current $T$ of
bidegree $(0,0)$ on $C^{\mathrm{an}}$ is of the form
\begin{equation}
                         T=c\,\delta_{C^{\mathrm{an}}},
            \qquad c\in\mathbf R_{\geq0}.
\label{eq:p213-triv-curve-classification}
\end{equation}
Thus every such $T$ is a Bedford--Taylor current in the full sense of
\cite[Definition~8.6.2]{CLD25}.

The same conclusion holds on the analytification of a connected Zariski
open subset of $C$.  The assertion is componentwise for a disjoint union.
\end{proposition}

\begin{proposition}
\label{prop:p213-triv-four-puncture-no-go}
Let $k$ be algebraically closed and trivially valued.  Choose pairwise
distinct $a,b,c,d\in k$, and put
\[
 X=\bigl(\mathbf P^1_k\setminus\{a,b,c,d\}\bigr)^{\mathrm{an}},
 \qquad
 f=\frac{z-a}{z-b},\qquad g=\frac{z-c}{z-d}.
\]
For
\[
 \tau=(-\log|f|,-\log|g|)=(s,t):X\longrightarrow\mathbf R^2,
\]
the image is the full cross with rays $e_1,-e_1,e_2,-e_2$.  The functions
\begin{equation}
                 U=2|s|,
       \qquad    V=|s|+|t|
\label{eq:p213-triv-cross-potentials}
\end{equation}
are continuous CLD-psh functions and satisfy
\begin{equation}
                 \mathrm d'\mathrm d''[U]
                  =\mathrm d'\mathrm d''[V]
                  =4\delta_\eta,
\label{eq:p213-triv-cross-curvature}
\end{equation}
where $\eta$ is the trivial absolute value on $k(z)$.

There is no positive closed current of bidegree $(0,0)$ whose branch
weights at $a,b,c,d$ are $2,2,1,1$.  More precisely, the four weights of
any positive closed current are equal.  Hence the natural weighted-cross
candidate cannot violate Bedford--Taylor independence: its visible
Bedford--Taylor defect is exactly the obstruction to closedness of the
putative multiplier current.
\end{proposition}

\begin{remark}
\label{rem:p213-triv-quarantine-scope}
The preceding direct calculations do not classify arbitrary
normal or nonnormal analytic curves over a trivially valued field, including
nonalgebraizable or singular curves.  For a smooth projective curve of
positive genus, the assertion above covers the projective analytification
and connected Zariski opens, but deliberately does not claim arbitrary
connected analytic opens.  The broader arbitrary-valued normal and
componentwise results are
\cref{prop:p213-normal-curve-arbitrary-valuation-current-revised,%
prop:p213-componentwise-arbitrary-valuation}.
\end{remark}
